\subsection*{1. Diversos casos particulares sobre un topos cualquiera}

1.1. En esta sección y en la siguiente, como en el Exposé VII, $T$
denota un topos, y $P,Q,G$ denotan tres grupos abelianos de $T$. Se
tiene, pues, un isomorfismo canónico (VII 3.6.5)
\begin{equation}\tag{1.1.1}
 \Biext^1(P,Q;G)
 \simeq
 \Ext^1\!\left(P\mathbin{\overset{\mathbb L}{\otimes}}Q,G\right).
\end{equation}
Recordemos los isomorfismos canónicos (llamados isomorfismos de Cartan)
\begin{equation}\tag{1.1.2}
\begin{split}
 \Ext^1\!\left(P\mathbin{\overset{\mathbb L}{\otimes}}Q,G\right)
 &\simeq \Ext^1\!\left(P,R\uHom(Q,G)\right)\\
 &\simeq \Ext^1\!\left(Q,R\uHom(P,G)\right).
\end{split}
\end{equation}
Todos estos isomorfismos son funtoriales en los tres argumentos
$P,Q,G$.

\sisi{Teniendo en cuenta que los objetos de homología de
$P\mathbin{\overset{\mathbb L}{\otimes}}Q$ son los
$\Tor_i(P,Q)$, que se anulan para $i\geq1$, se deduce de (1.1.1) la
siguiente sucesión exacta infinita:}{Tomando el truncamiento estándar
del producto tensorial derivado, cuyos objetos de homología son los
$\Tor_i(P,Q)$, se obtiene de (1.1.1) la siguiente sucesión exacta
infinita:}
\begin{equation}\tag{1.1.3}
\begin{aligned}
0\longrightarrow{}&
 \Ext^1(P\otimes Q,G)
 \longrightarrow \Biext^1(P,Q;G)
 \longrightarrow \Hom(\Tor_1(P,Q),G)\\
 \longrightarrow{}&
 \Ext^2(P\otimes Q,G)
 \longrightarrow
 \Ext^2\!\left(P\mathbin{\overset{\mathbb L}{\otimes}}Q,G\right)
 \longrightarrow\cdots .
\end{aligned}
\end{equation}
Por otra parte, como los objetos de cohomología de $R\uHom(Q,G)$ son
los $\uExt^i(Q,G)$, la primera igualdad de (1.1.2) da lugar a la
sucesión exacta de cinco términos
\begin{equation}\tag{1.1.4}
\begin{aligned}
0\longrightarrow{}&
 \Ext^1(P,\uHom(Q,G))
 \longrightarrow \Biext^1(P,Q;G)
 \longrightarrow \Hom(P,\uExt^1(Q,G))\\
 \longrightarrow{}&
 \Ext^2(P,\uHom(Q,\sisi{P}{G}))
 \longrightarrow \Ext^2(P,R\uHom(Q,G)),
\end{aligned}
\end{equation}

%%%%%%%%%%  scan idx 232  |  folio 221  %%%%%%%%%%

asociada a la sucesión espectral usual
\[
 \Ext^*(P,R\uHom(Q,G))
 \Longleftarrow
 \Ext^p(P,\uExt^q(Q,G)).
\]
Dejamos al lector, si siente la necesidad, la tarea de explicitar la
primera flecha de (1.1.3), respectivamente de (1.1.4), por medio del
\underline{funtor natural plenamente fiel}
\begin{equation}\tag{1.1.5}
 \operatorname{EXT}(P\otimes Q,G)
 \longrightarrow
 \operatorname{BIEXT}(P,Q;G),
\end{equation}
respectivamente
\begin{equation}\tag{1.1.6}
 \operatorname{EXT}(P,\uHom(Q,G))
 \longrightarrow
 \operatorname{BIEXT}(P,Q;G),
\end{equation}
donde $\operatorname{EXT}$ y $\operatorname{BIEXT}$ denotan las
categorías evidentes de extensiones y biextensiones.

Naturalmente, el isomorfismo
\[
 \Biext^1(P,Q;G)\simeq\Ext^1(Q,R\uHom(P,G))
\]
da lugar a una sucesión exacta análoga a (1.1.3), con los papeles de
$P$ y $Q$ intercambiados. Su primera flecha puede explicitarse por
medio del funtor plenamente fiel
\begin{equation}\tag{1.1.7}
 \operatorname{EXT}(Q,\uHom(P,G))
 \longrightarrow
 \operatorname{BIEXT}(P,Q;G)
\end{equation}
análogo a (1.1.6).

1.2. Supongamos que
\begin{equation}\tag{1.2.1}
 \Tor_1(P,Q)=0.
\end{equation}
Entonces (1.1.3) proporciona un isomorfismo canónico
\begin{equation}\tag{1.2.2}
 \Biext^1(P,Q;G)\simeq\Ext^1(P\otimes Q,G).
\end{equation}
Así ocurre, por ejemplo, si $P$ o $Q$ es un $\mathbb Z$-módulo
\underline{plano}. Más precisamente, bajo la condición (1.2.1), el
funtor (1.1.5) es una equivalencia de categorías.

%%%%%%%%%%  scan idx 233  |  folio 222  %%%%%%%%%%

1.3. Supongamos que
\begin{equation}\tag{1.3.1}
 P\otimes Q=0.
\end{equation}
Entonces (1.1.3) proporciona un isomorfismo canónico
\begin{equation}\tag{1.3.2}
 \Biext^1(P,Q;G)\simeq\Hom(\Tor_1(P,Q),G).
\end{equation}
El caso de mayor interés para nosotros en el que se cumple (1.3.1) es
aquel en que existe un entero $n>0$ tal que
\begin{equation}\tag{1.3.3}
 nP=P,\qquad
 Q=\varinjlim_i {}_{n^i}Q;
\end{equation}
es decir, $P$ es divisible por $n$ y $Q$ es límite de subgrupos
anulados por potencias de $n$. Pongamos, como de costumbre,
\[
 T_n(P)=({}_{n^i}P)_{i\geq0}\quad\text{(un sistema proyectivo)},\qquad
 {}_{n^\infty}P=\varinjlim_i{}_{n^i}P.
\]
De (1.3.3) se obtienen las fórmulas siguientes, más precisas que
(1.3.1):
\begin{equation}\tag{1.3.4}
\begin{split}
 P\otimes Q&=0,\\
 \Tor_1(P,Q)
 &\simeq\varinjlim_i
   \bigl({}_{n^i}P\otimes{}_{n^i}Q\bigr)
 \simeq T_n(P)\otimes Q\\
 &\simeq{}_{n^\infty}P\otimes T_n(Q).
\end{split}
\end{equation}
Por tanto, el isomorfismo (1.3.2) adopta aquí la forma más concreta
\begin{equation}\tag{1.3.5}
\begin{split}
 \Biext^1(P,Q;G)
 &\simeq\Hom(T_n(P),D(Q))\\
 &\simeq\Hom(Q,D(T_n(P))\sisi{}{)}\\
 &\simeq\Hom(T_n(P)\otimes T_n(Q),T_n(G)).
\end{split}
\end{equation}
Precisemos el significado de estas fórmulas.

\begin{enumerate}[label=\alph*)]
\item En el sistema inductivo que aparece en (1.3.4), el homomorfismo
de transición
\[
 {}_{n^i}P\otimes{}_{n^i}Q
 \longrightarrow
 {}_{n^j}P\otimes{}_{n^j}Q
 \qquad(i\leq j)
\]
se describe identificando el primer miembro con
$\,{}_{n^j}P\otimes{}_{n^i}Q$ mediante
$n^{j-i}\otimes\id$, y utilizando después el homomorfismo
\[
 \id\otimes\operatorname{inj}:
 {}_{n^j}P\otimes{}_{n^i}Q
 \longrightarrow
 {}_{n^j}P\otimes{}_{n^j}Q,
\]
donde
$\operatorname{inj}:{}_{n^i}Q\longrightarrow{}_{n^j}Q$
denota la inclusión.

%%%%%%%%%%  scan idx 234  |  folio 223  %%%%%%%%%%

\item Por definición, pongamos
\[
 T_n(P)\otimes Q
 \overset{\mathrm{def}}{=}
 \varinjlim_i\bigl(T_n(P)\otimes{}_{n^i}Q\bigr),
\]
donde
\[
 T_n(P)\otimes M
 \overset{\mathrm{def}}{=}
 {}_{n^i}P\otimes M
 \qquad\text{si }n^iM=0.
\]
El tercer término escrito en (1.3.4) se precisa del mismo modo cuando
se supone también $nQ=Q$, de manera que $T_n(Q)$, como $T_n(P)$, es
un sistema proyectivo estricto. El último término de
\sisi{(1.3.4)}{(1.3.5)} es un $\Hom$ entre sistemas proyectivos de
grupos.

\item La letra $D$ denota $\uHom(-,G)$. El objeto $D(Q)$ se define
como el sistema proyectivo de grupos
\[
 D(Q)=\bigl(D({}_{n^i}Q)\bigr)_{i\geq0}
 =\bigl(\uHom({}_{n^i}Q,G)\bigr)_{i\geq0},
\]
siendo el segundo miembro de (1.3.5) el grupo de homomorfismos de
sistemas proyectivos de grupos. \sisi{El último miembro}{El miembro
siguiente} se interpreta más sencillamente poniendo
\[
 D(T_n(P))=\varinjlim_iD({}_{n^i}P),
\]
donde los morfismos de transición de este sistema inductivo se
obtienen por transposición de los morfismos de transición
\[
 \sisi{n^{j-1}}{n^{j-i}}:{}_{n^j}P\longrightarrow{}_{n^i}P
\]
de $T_n(P)$.
\end{enumerate}

%%%%%%%%%%  scan idx 235  |  folio 224  %%%%%%%%%%

Dejamos al lector la tarea de explicitar el cálculo de $\Tor_1(P,Q)$
indicado en (1.3.4) y de deducir de él (1.3.5). Señalemos únicamente
que se utilizan las relaciones, deducidas de la segunda hipótesis de
(1.3.3),
\[
 \Tor_p(P,Q)\simeq
 \varinjlim_i\Tor_p(P,{}_{n^i}Q).
\]
Para $p=0$ esto da inmediatamente $P\otimes Q=0$ y reduce el cálculo
de $\Tor_1(P,Q)$ al de los $\Tor_1(P,{}_{n^i}Q)$ y sus morfismos de
transición. Así, cuando $Q$ está anulado por un entero $m=n^i$ tal que
$mP=P$, queda por explicitar un isomorfismo
\begin{equation}\tag{1.3.6}
 \Tor_1(P,Q)
 \xleftarrow{\ \sim\ }
 {}_mP\otimes{}_mQ
 \qquad(mP=P,\ mQ=0).
\end{equation}
Hay varias maneras de hacerlo. Una consiste en utilizar la sucesión
exacta de Tor asociada a
\[
 0\longrightarrow{}_mP\longrightarrow P
 \xrightarrow{\ m\ }P\longrightarrow0.
\]
También puede utilizarse el homomorfismo (2.1.1) de más abajo; el hecho
de que los morfismos de transición sean los precisados en (a) no es
sino 2.1.11. En cualquier caso, la elección de (1.3.6) equivale a la
elección de un signo; véase 2.1.10.

1.4. Supongamos que
\begin{equation}\tag{1.4.1}
 \uHom(Q,G)=0.
\end{equation}
Entonces (1.1.4) proporciona un isomorfismo canónico
\begin{equation}\tag{1.4.2}
 \Biext^1(P,Q;G)
 \xrightarrow{\ \sim\ }
 \Hom(P,\uExt^1(Q,G)).
\end{equation}

%%%%%%%%%%  scan idx 236  |  folio 225  %%%%%%%%%%

Este caso se presenta en particular cuando $Q$ es un esquema abeliano
y $G=\underline G_m$; véase VII 1.3.8. Por consiguiente, el funtor en
$P$,
\[
 P\longmapsto\Biext^1(P,Q;G),
\]
es \underline{representable} por $Q^*=\uExt^1(Q,G)$. En particular,
existe una biextensión canónica $E^Q$ de $(Q^*,Q)$ por $G$,
determinada salvo isomorfismo único, pues su grupo de automorfismos es
\[
 \Hom(\sisi{P}{Q^*},\uHom(Q,G))=0.
\]
Esta biextensión, o su biextensión simétrica $\,{}^sE^Q$ de
$(Q,Q^*)$ por $G$, puede llamarse la \underline{biextensión de Weil}
de $(Q^*,Q)$, respectivamente de $(Q,Q^*)$, por $G$.

1.5. Supongamos que
\begin{equation}\tag{1.5.1}
 \Hom(P,\uExt^1(Q,G))=0.
\end{equation}
Esto se cumple en particular si $\uExt^1(Q,G)=0$; por ejemplo, si $Q$
es un esquema en grupos finito localmente libre o un grupo de tipo
multiplicativo (SGA 3 IX), y $G=\underline G_m$; véase 3.4 más abajo.
Entonces (1.1.4) proporciona un isomorfismo canónico
\begin{equation}\tag{1.5.2}
 \Biext^1(P,Q;G)
 \xleftarrow{\ \sim\ }
 \Ext^1(P,\uHom(Q,G)),
\end{equation}
donde $\uHom(Q,G)$ desempeña el papel de un ``grupo dual'' de $Q$
(relativamente a $G$). Más precisamente, de (1.5.2) se sigue que el
funtor plenamente fiel (1.1.6) es una \underline{equivalencia de
categorías}.

1.6. Supongamos que se cumplen simultáneamente las dos relaciones
\begin{equation}\tag{1.6.1}
 \uHom(P,G)=0,\qquad
 \Hom(P,\uExt^1(Q,G))=0.
\end{equation}
Por 1.4 y 1.5, respectivamente, estas dos relaciones dan los
isomorfismos
\[
 \Biext^1(P,Q;G)
 \xrightarrow{\ \sim\ }
 \Hom(Q,\uExt^1(P,G)),
\]
\[
 \Biext^1(P,Q;G)
 \xleftarrow{\ \sim\ }
 \Ext^1(P,\uHom(Q,G)).
\]
Su compuesto es el isomorfismo canónico
\begin{equation}\tag{1.6.2}
 \Ext^1(P,\uHom(Q,G))
 \xrightarrow{\ \sim\ }
 \Hom(Q,\uExt^1(P,G))
 \quad\bigl(\simeq\Biext^1(P,\sisi{G}{Q};G)\bigr),
\end{equation}

%%%%%%%%%%  scan idx 237  |  folio 226  %%%%%%%%%%

que desempeña el papel de una fórmula de dualidad.

Obsérvese además que la primera relación de (1.6.1) implica
evidentemente $\Hom(P\otimes Q,G)=0$. Por tanto, la categoría
$\operatorname{EXT}(P,\uHom(Q,G))$ de extensiones conmutativas de $P$
por $\uHom(Q,G)$ es rígida (todo automorfismo de un objeto es la
identidad), y lo mismo ocurre con $\operatorname{BIEXT}(P,Q;G)$. La
estructura de estas categorías queda, pues, determinada cuando se
conoce el grupo de clases de isomorfismo de sus objetos; por (1.6.2),
este grupo es canónicamente isomorfo a
\[
 \Hom(Q,\uExt^1(P,\sisi{Q}{G})).
\]

La aplicación (1.6.2) está de hecho definida en todos los casos, sin
suponer (1.6.1), y también puede describirse directamente, sin pasar
por $\Biext^1(P,Q;G)$, en términos del acoplamiento natural
\[
 Q\times Q'\longrightarrow G,
 \qquad Q'=\uHom(Q,G).
\]
Interpretándolo como un homomorfismo
$Q\longrightarrow\uHom(Q',G)$ se obtiene un compuesto
\begin{equation}\tag{1.6.3}
\begin{split}
 \Ext^1(P,Q')\times Q
 &\longrightarrow
 \Ext^1(P,Q')\times\Hom(Q',G)\\
 &\longrightarrow
 \Ext^1(P,G),
\end{split}
\end{equation}
que vuelve a dar (1.6.2).

Para un caso útil en el que se aplican las consideraciones de 1.6,
véase 3.7 más abajo.
